\documentclass[10pt,a4paper]{amsart}
\usepackage[margin=19mm]{geometry}
\usepackage{amsmath,amssymb,mathtools}
\usepackage[scr=rsfs,cal=euler]{mathalfa}
\usepackage{xfrac}
\usepackage[unicode,hidelinks,bookmarks=false]{hyperref}
\newcommand{\ie}{i.e.\ }
\newcommand{\cf}{cf.\ }
\newcommand{\resp}{respectively\ }
\newcommand{\Subsection}{\subsection}
\providecommand{\nobreakdash}{\nobreak\hbox{-}}
\setlength{\emergencystretch}{2em}
\multlinegap0pt
\usepackage{float}
\usepackage{csquotes}
\usepackage{amssymb}
\newtheorem{lemma}{Lemma}
\title{Interpolations: a linear algebra approach}
\author{Andronick Arutyunov, Nikolay Korgin}
\date{Source: arXiv:2606.22671v1 (CC BY 4.0)}
\begin{document}
\maketitle

\noindent\emph{Source and licence.} Interpolations: a linear algebra approach. Version arXiv:2606.22671v1 (CC BY 4.0), distributed by arXiv under the Creative Commons Attribution 4.0 International licence, http://creativecommons.org/licenses/by/4.0/, as recorded in the arXiv licence metadata for this exact version and shown on the abstract page https://arxiv.org/abs/2606.22671v1. Full licence notice: \texttt{licenses/arxiv-2606.22671-CC-BY-4.0.txt}.\par\smallskip

\noindent\textit{Editorial scope.} Counted unit: the article's lemma lemma-V+V** (source0.tex lines 402--404). The article's complete original proof of it is delivered with it (source0.tex lines 406--436, one \texttt{proof} environment of 31 lines). No mathematics is written, repaired or re-derived: the statement and the proof are the article's own lines, in the article's own notation, and no retained line is substituted or re-flowed.  No further source-local context is needed: every cross-reference of every retained line resolves inside the two delivered spans.  Nothing was omitted from any retained span, so the package carries no omission record.  No retained line contains a citation, so the wrapper carries no bibliography and no bibliography entry.  No retained line contains a figure environment, an image-inclusion command or drawing code, so no image is delivered and the package declares no asset dependency.  Editorial formatting only, in addition: an AMS \texttt{amsart} preamble that restates the article's own declarations and macros used by the retained lines, namely the environments \texttt{lemma} on separate counters, together with the standard packages those lines need.  The retained environments print in this document as Lemma 1 in order of appearance, whereas the article prints the same items with the numbering of its own full document.  The complete native source of the article as distributed in the arXiv e-print for this exact version is bundled unchanged as \texttt{sources/203372/source0.tex}.
\begin{lemma}\label{lemma-V+V**}
If $V = V_1 \oplus V_2$, then the dual basis corresponding to the union of bases of $V_1$ and $V_2$ is the union of the dual bases of $V_1^*$ and $V_2^*$.
\end{lemma}

\begin{proof}
We first show that
\[
    V^* = V_1^* \oplus V_2^*.
\]
To this end, extend any linear functional $\varphi \in V_1^*$ to a linear functional on $V$ in a trivial way: for every $v \in V_2$, set $\varphi(v) = 0$.

Now construct the dual basis for the decomposed spaces. Suppose we have bases
\[
    V_1 = \langle e^1_i \rangle,\quad i = 1..n_1,
    \qquad
    V_2 = \langle e^2_j \rangle,\quad j = 1..n_2,
\]
and the corresponding dual bases
\[
    V_1^* = \langle \varepsilon^1_i \rangle,\quad i = 1..n_1,
    \qquad
    V_2^* = \langle \varepsilon^2_j \rangle,\quad j = 1..n_2.
\]

Since $\varepsilon^1_i(v_2) = 0$ for every $v_2 \in V_2$, and $\varepsilon^2_j(v_1) = 0$ for every $v_1 \in V_1$, it follows that the set
\[
    V^* = \langle \varepsilon^1_i,\, \varepsilon^2_j \rangle,
    \qquad i = 1..n_1,\; j = 1..n_2,
\]
is the dual basis corresponding to the basis
\[
    V = \langle e^1_i,\, e^2_j \rangle,
    \qquad i = 1..n_1,\; j = 1..n_2.
\]
\end{proof}

\end{document}
